% Generated draft: requires mathematical and source review.
\documentclass[11pt]{article}
\usepackage[margin=25mm]{geometry}
\usepackage{fontspec}
\setmainfont{Latin Modern Roman}
\usepackage{amsmath,amssymb,mathtools,unicode-math}
\setmathfont{Latin Modern Math}
\usepackage{enumitem,needspace,xurl,hyperref}
\setlength{\parindent}{0pt}
\setlength{\parskip}{5pt}
\setlength{\emergencystretch}{3em}
\newcommand{\N}{\mathbb{N}}
\newcommand{\Z}{\mathbb{Z}}
\newcommand{\Q}{\mathbb{Q}}
\newcommand{\R}{\mathbb{R}}
\newcommand{\C}{\mathbb{C}}
\newcommand{\F}{\mathbb{F}}
\newcommand{\A}{\mathbb{A}}
\newcommand{\PP}{\mathbb{P}}
\newcommand{\E}{\mathbb{E}}
\newcommand{\eps}{\varepsilon}
\newcommand{\dd}{\,\mathrm{d}}
\begin{document}
\section*{UnsolvedMath ID: 30003079}
\textbf{Source problem code:} catalogue entry
\subsection{Definitions and mathematical statement}
The supplied question is: ``Is Terao's conjecture true for line arrangements; equivalently, is freeness determined by the incidence lattice?'' The source does not specify the ground field, its characteristic, or whether the lines are affine or projective. These choices remain unresolved. The following gives a precise projective formulation conditional on fixing a ground field \(k\); it does not assert that the source quantifies over all fields.

Let \(\mathcal A\) be a finite set of distinct projective lines in \(\PP^2_k\). Choose a nonzero homogeneous linear equation \(\alpha_H\in S=k[x_0,x_1,x_2]\) for each \(H\in\mathcal A\). Its central cone is the arrangement \(c\mathcal A=\{\ker(\alpha_H)\subseteq k^3:H\in\mathcal A\}\). Define its intersection lattice by
\[
L(c\mathcal A)=\left\{\bigcap_{H\in\mathcal B}\ker(\alpha_H):\mathcal B\subseteq\mathcal A\right\},
\]
where the empty intersection is \(k^3\), and order this set by reverse inclusion. This is the convention for the incidence lattice used in this conditional formulation.

The module of logarithmic derivations is
\[
D(c\mathcal A)=\left\{\theta\in\operatorname{Der}_k(S):\theta(\alpha_H)\in\alpha_H S\ \text{for every }H\in\mathcal A\right\}.
\]
Here \(\operatorname{Der}_k(S)\) denotes the \(S\)-module of \(k\)-linear derivations of \(S\). Call \(\mathcal A\) free when \(D(c\mathcal A)\) is a free \(S\)-module. These definitions are independent of the nonzero scalar choices of the equations \(\alpha_H\).

For the fixed field \(k\), the projective version of the question asks whether, for every pair \(\mathcal A,\mathcal B\) of such arrangements over that same field,
\[
L(c\mathcal A)\cong L(c\mathcal B)
\quad\Longrightarrow\quad
\bigl(D(c\mathcal A)\text{ is free over }S\iff D(c\mathcal B)\text{ is free over }S\bigr).
\]
The intended field and geometric convention must be confirmed before identifying this conditional formulation with the full intended source statement.
\subsection{Short English statement}
The question asks whether the intersection data of a line arrangement determine its freeness. A precise projective version compares arrangements over one fixed field and defines freeness through logarithmic derivations of their central cones. The supplied record leaves the field and the affine or projective convention unspecified.
\subsection{Sources}
\begin{itemize}
\item Local UnsolvedMath registry, permanent ID \(30003079\), title ``Terao Conjecture for Line Arrangements,'' source problem code \texttt{catalogue entry}, in the Oberwolfach Reports Open Problems collection. The record associates the problem with 2016 and supplies no individual proposer. It labels the statement ``exact,'' with the supplied verification that it is internally readable and complete and was retained byte-for-byte; this is not independent mathematical verification.
\item Supplied primary provenance: \emph{Mini-Workshop: Arrangements of Subvarieties, and their Applications in Algebraic Geometry} (2016), Oberwolfach Reports, pp. 653--702, DOI 10.4171/owr/2016/14. \url{https://doi.org/10.4171/owr/2016/14}. The supplied background describes Oberwolfach Reports as documenting mathematical workshops and their open-problem discussions.
\item Supplied literature source: ``Pagaria.'' \url{https://doi.org/10.1112/jlms.12829}.
\item Supplied literature source: ``Pokora et al.'' \url{https://doi.org/10.1007/s40879-025-00873-2}.
\item Literature assessment dated exactly \texttt{2026-08-21}: the supplied curation labels the problem open and states that the report identifies deciding Terao's conjecture as a main open problem, with later work establishing additional classes but not all line arrangements. This is supplied dated curation derived from the OWR record, not an independently verified assessment. Its time-sensitive status and cited claims require independent verification.
\end{itemize}
\end{document}
